% =========================
% 8_appendix.tex --- follows the bibliography; does not count towards the 12-page limit.
% Nothing load-bearing for the argument lives here: reviewers are not obligated to read it.
% =========================
\appendix

\section{The Dependency Structure Behind the Combination}
\label{app:entailment}

Section~\ref{sec:combination} claims HALO's four choices are \emph{mutually entailed} rather than
independently attractive. Table~\ref{tab:coverage} records which axes the closest systems claim to
address; the four dependencies follow in full.

\begin{table}[h]
  \centering
  \caption{Which axes each system \emph{explicitly claims} to address. $\checkmark$~= in the form
  we argue is needed; $\circ$~= a weaker form (rate \emph{equivariance} rather than physical-Hz
  anchoring; \emph{categorical} rather than free-text placement; rotation \emph{invariance},
  which discards gravity). The last column marks systems reporting a range of label budgets with
  no gradient update. Blank is not a criticism---these are different papers solving different
  problems---and the table reports stated scope, not measured performance. The point is not the
  bottom row's length but the dependency structure below.}
  \label{tab:coverage}
  \scriptsize
  \setlength{\tabcolsep}{2.5pt}
  \begin{tabular}{@{}l ccccc@{}}
    \toprule
    & \multicolumn{2}{c}{acquisition} & orient./ & label & $k$-curve \\
    \cmidrule(lr){2-3}
    System & rate & place & grav. & interf. & \\
    \midrule
    LiMU-BERT~\cite{limubert}, CrossHAR~\cite{crosshar}\hspace*{-4pt} &   &   &   &   &  \\
    harnet~\cite{sslwearables}, LSM~\cite{lsm1,lsm2}    &   &   &   &   & \checkmark \\
    NormWear~\cite{normwear}                &         & $\circ$ &    & \checkmark &  \\
    UniMTS~\cite{unimts}                    &         & \checkmark & $\circ$ & \checkmark &  \\
    ZeroHAR~\cite{zerohar}                  &         & $\circ$ &        & \checkmark &  \\
    GOAT~\cite{goat}, oneHAR~\cite{onehar}  &         & \checkmark &     & \checkmark &  \\
    CHARM~\cite{charm,sensorsvoice}         &         & \checkmark &     &        &  \\
    MoPFormer~\cite{mopformer}              &         & $\circ$ &        &        &  \\
    FlowState~\cite{flowstate}              & $\circ$ &         &        &        &  \\
    RAG-HAR~\cite{raghar}, ZARA~\cite{zara} &         &         &        & \checkmark & \checkmark \\
    SensorLM~\cite{sensorlm}                &         &         &        & \checkmark & \checkmark \\
    \midrule
    \textbf{HALO}                           & \checkmark & \checkmark & \checkmark & \checkmark & \checkmark \\
    \bottomrule
  \end{tabular}
\end{table}

\textbf{Rate anchoring is a precondition for the memory, not a parallel feature.} A
non-parametric memory only works if a query and a stored exemplar recorded on different devices
are comparable \emph{as vectors}. If the encoder is rate-\emph{equivariant}---told the rate and
adapting---then two windows at different rates are related by a learned transformation, and the
memory must either be indexed per configuration or trust that transformation to be exact. Because
HALO's band $k$ denotes the same physical frequency for every device, one flat memory is
well-posed. L1 exists in the form it does because of Phase~B.

\textbf{The granularity of the language conditioning is forced by the retrieval.} Per-channel
free text~\cite{charm,sensorsvoice,sensorllm} is the natural choice if the conditioning feeds a
classifier: it is maximally expressive and the head can ignore what it does not need. It is the
wrong choice if the conditioning feeds a memory, because a description nearly unique to a corpus
makes \emph{same-corpus} neighbours the nearest ones---precisely the shortcut cross-dataset
evaluation exists to catch. Attaching identity to the \emph{sensor} keeps the description
expressive enough to name an unseen device and too coarse to name a dataset. That is why the
dataset-identity probe is in the paper.

\textbf{Preserving gravity is forced by the geometry finding.} Rotation
invariance~\cite{unimts} is the standard answer to unknown mounting, and it is lossy in one
specific way: it discards the gravity direction, and with it posture. Section~\ref{sec:geometry}
shows that the label pairs a text projection cannot separate are exactly the posture-versus-
transition pairs (\emph{sitting}/\emph{sitting down}, \emph{lying}/\emph{standing}). A system
that gives up on the text geometry \emph{must} therefore keep posture in the signal, so
conditioning rather than invariance is not a preference here but a requirement.

\textbf{Label-free pretraining is forced by the same shortcut argument.} Any closed-vocabulary
objective over a multi-corpus training set rewards inferring the corpus and predicting its
marginal. That is survivable for a parametric head, which can be refitted per deployment; it is
not survivable for a frozen encoder feeding a memory, because the shortcut is baked into the
stored vectors.

So the four choices form a chain: removing any one invalidates another. That is a different and
more falsifiable claim than ``we do all of these,'' and Section~\ref{sec:bias_experiments} is
built to test it---each ablation in Table~\ref{tab:tokenizer_ablation} removes one link and asks
whether the rest still stand.

\section{Corpus}
\label{app:corpus}

Table~\ref{tab:datasets} in the body lists the training corpus and the held-out evaluation sets;
Figure~\ref{fig:corpus} plots them by placement and native rate. Window counts there are windows
on the native-rate grid after physical-plausibility and duplicate filtering. Most windows are six
seconds, but UCI-HAR ships $2.56$\,s records and SP-SW-HAR $1$\,s, so hours are computed from each
grid's own samples and rate rather than by multiplying by six.

\begin{figure}[h]
    \centering
    \setlength{\abovecaptionskip}{2pt}
  % =====================================================================
% fig_corpus.tex  --  corpus heterogeneity: placement x sampling rate
%
% DATA PROVENANCE: every point is a real stream in the materialized native
% grids, enumerated 2026-07-26 from training.tokenizer.pretrain_data.CorpusIndex
% (20 training streams over 12 datasets) and data/datasets/*/metadata.json
% (7 held-out datasets).  Bubble area encodes window count; the largest
% (CAPTURE-24 wrist) is 144,142 windows, the smallest (SP-SW-HAR pocket) 1,179.
% =====================================================================
\resizebox{\columnwidth}{!}{%
\begin{tikzpicture}
\begin{axis}[
  width=\columnwidth, height=4.6cm,
  xmin=0.55, xmax=5.55,
  ymin=6, ymax=145,
  ytick={20,50,100},
  yticklabels={20\,Hz,50\,Hz,100\,Hz},
  yticklabel style={font=\scriptsize},
  xtick={1,2,3,4,5},
  xticklabels={waist/hip,pocket,wrist,lower back,head/ear},
  xticklabel style={font=\scriptsize},
  ylabel={sampling rate},
  ylabel style={font=\scriptsize,yshift=-6pt},
  grid=major, grid style={draw=halogray!18},
  axis line style={draw=halogray!60},
  tick style={draw=halogray!60},
  scale only axis,
  clip=false,
]
% ---------- training streams (filled blue discs, diameter ~ #windows^0.4) ----
% diameters below are 2*(1.2 + 4.5*(n/144142)^0.4) pt, n = window count
\tikzset{hdisc/.style={circle,draw=haloblue,fill=haloblue!45,line width=0.4pt,inner sep=0pt}}
\node[hdisc,minimum size=5.44pt] at (axis cs:0.88,50)  {};% hhar        9,587
\node[hdisc,minimum size=5.40pt] at (axis cs:1.00,100) {};% kuhar       9,233
\node[hdisc,minimum size=5.54pt] at (axis cs:1.12,50)  {};% uci_har    10,299
\node[hdisc,minimum size=5.70pt] at (axis cs:1.80,50)  {};% unimib     11,771
\node[hdisc,minimum size=4.88pt] at (axis cs:1.94,50)  {};% xrf l.pkt   5,794
\node[hdisc,minimum size=7.76pt] at (axis cs:2.00,20)  {};% wisdm pkt  39,329
\node[hdisc,minimum size=4.88pt] at (axis cs:2.06,50)  {};% xrf r.pkt   5,794
\node[hdisc,minimum size=3.72pt] at (axis cs:2.00,100) {};% spsw pkt    1,179
\node[hdisc,minimum size=11.4pt] at (axis cs:2.82,100) {};% capture24 144,142
\node[hdisc,minimum size=6.36pt] at (axis cs:2.85,50)  {};% harmes     18,576
\node[hdisc,minimum size=4.36pt] at (axis cs:2.94,100) {};% pamap2      3,186
\node[hdisc,minimum size=3.68pt] at (axis cs:2.95,50)  {};% mhealth     1,104
\node[hdisc,minimum size=7.26pt] at (axis cs:3.00,20)  {};% wisdm wr   30,876
\node[hdisc,minimum size=4.88pt] at (axis cs:3.05,50)  {};% xrf l.wr    5,794
\node[hdisc,minimum size=5.86pt] at (axis cs:3.06,100) {};% nfi wrist  13,260
\node[hdisc,minimum size=4.88pt] at (axis cs:3.15,50)  {};% xrf r.wr    5,794
\node[hdisc,minimum size=4.20pt] at (axis cs:3.18,100) {};% spsw wrist  2,578
\node[hdisc,minimum size=5.86pt] at (axis cs:4.00,100) {};% nfi back   13,260
\node[hdisc,minimum size=4.88pt] at (axis cs:4.92,50)  {};% xrf glasses 5,794
\node[hdisc,minimum size=4.88pt] at (axis cs:5.08,50)  {};% xrf earbuds 5,794
% ---------- held-out datasets (open orange triangles) ----------
\addplot[only marks,mark=triangle*,mark size=2.6pt,color=haloorange,
         fill=halolight2,line width=0.7pt]
  coordinates {
    (1.26,50) (1.38,50) (1.26,100)
    (2.20,50) (2.32,50)
    (3.30,50) (3.42,50)
  };
% ---------- annotations ----------
\node[font=\scriptsize,color=haloblue!75!black,anchor=west] at (axis cs:2.95,128)
  {\textbf{CAPTURE-24} 144k windows};
\draw[haloblue!70,thin] (axis cs:2.86,106) -- (axis cs:2.93,124);
\node[font=\scriptsize,color=halogray,anchor=west,align=left] at (axis cs:4.05,24)
  {smart glasses\\ \& earbuds};
\draw[halogray,thin] (axis cs:5.00,44) -- (axis cs:4.90,31);
\node[font=\scriptsize,anchor=west,align=left] at (axis cs:0.62,13)
  {\textcolor{haloblue!75!black}{$\bullet$ 20 training streams}\;\;
   \textcolor{haloorange!85!black}{$\blacktriangle$ 7 held-out}};
\end{axis}
\end{tikzpicture}}
\caption{\textbf{The corpus does not cover a grid; it covers a scatter.} Each blue disc is
one of the 20 sensor streams HALO trains on, positioned by body placement and native
sampling rate, with area proportional to the number of $6$\,s windows it contributes; orange
triangles are the 7 held-out evaluation datasets. Rates span a $5\times$ range and placements
include two---smart glasses and wireless earbuds---that no established HAR benchmark
contains. Held-out cells are not simply new subjects: they combine placements and rates in
proportions the training corpus never sees. Absorbing this by resampling to a common rate
would discard the content above $10$\,Hz that only the $50$/$100$\,Hz streams can
resolve.}
\label{fig:corpus}

\end{figure}

\begin{figure*}[t]
    \centering
    \setlength{\abovecaptionskip}{4pt}
    \setlength{\belowcaptionskip}{-6pt}
  \resizebox{0.92\textwidth}{!}{% Retired from the manuscript. The former teaser embedded trace values without
% a checked-in generator that bound them to source windows and artifact hashes.
% Reintroduce it only with provenance-complete generated data.
}
  \caption{\textbf{What ``no closed grammar'' looks like in the corpus.} The three traces on the
  left are real accelerometer snippets of \emph{walking}, two seconds each, from three streams of
  the materialised corpus; they are amplitude-normalised for display, but the sample counts are
  the true ones. As arrays they have nothing in common. HALO absorbs that heterogeneity
  \emph{in the tokenizer}---a frequency axis fixed in Hz, sensor identity supplied in language,
  gravity preserved---so windows recorded on different devices become directly comparable as
  vectors, which is the precondition for one flat memory. Recognition is then a retrieval in
  \emph{sensor} space, with the label strings a deployment declares at run time used only to name
  and to enumerate; because the memory is append-only, one frozen model spans the whole label
  budget with no gradient update at any point.}
  \label{fig:hook}
\end{figure*}

\textbf{Deliberate exclusions.} Six datasets were excluded from both splits.
DSADS~\cite{dsads}, HARTH~\cite{harth}, and Opportunity~\cite{opportunity} use torso and limb IMU
harnesses rather than phone or watch placements. RecGym~\cite{recgym} is distributed per-axis
min--max normalised, which destroys the physical scale and gravity component the tokenizer
depends on. MobiAct~\cite{mobiact} was never materialised into grids. And HAPT~\cite{hapt} is
excluded because it re-records the same 30 subjects as UCI-HAR~\cite{ucihar}, so including both
would place the same people on both sides of a supposedly subject-disjoint boundary.

\section{Baselines}
\label{app:baselines}

Table~\ref{tab:baselines} lists the five baselines, what each was pretrained on, and how much
acquisition context each can receive at inference---the axis on which the comparison is not
symmetric.

\begin{table}[h]
  \centering
  \caption{Baselines and where each sits on the heterogeneity stack. ``Blind'' means the model
  receives no information about where the sensor was worn. Two baselines are pretrained by us on
  our corpus with their published objectives (data-matched); three are used as released
  (faithful, but pretrained on their own, larger corpora). We report the two groups separately
  rather than averaging over them.}
  \label{tab:baselines}
  \scriptsize
  \setlength{\tabcolsep}{2.5pt}
  \begin{tabular}{@{}l l l l l@{}}
    \toprule
    Model & Weights & Rate & Channels / placement & Labels \\
    \midrule
    LiMU-BERT~\cite{limubert}   & we pretrain & 20\,Hz & 6-ch pad$+$mask; blind & ConSE \\
    CrossHAR~\cite{crosshar}    & we pretrain & 20\,Hz & 6-ch pad$+$mask; blind & ConSE \\
    harnet~\cite{sslwearables}  & released    & 30\,Hz & 3-ch acc; wrist & ConSE \\
    UniMTS~\cite{unimts}        & released    & resamp. & joint-graph placement & native \\
    NormWear~\cite{normwear}    & released    & 65\,Hz & variable ch; blind & native \\
    \midrule
    \textbf{HALO} (ours)        & we train    & \textbf{native} & \textbf{per-sensor language} & \textbf{native} \\
    \bottomrule
  \end{tabular}
\end{table}

\section{Parity Control}
\label{app:parity}

Because several baselines cannot accept HALO's native inputs, Table~\ref{tab:parity} reports HALO
restricted to each baseline's own input contract, so any remaining gap is attributable to the
model rather than to what it was allowed to read.

\begin{table}[h]
  \centering
  \caption{Parity control: HALO under its native inputs versus baseline-equivalent inputs.
  Baselines are split by whether their weights were trained on our corpus or released by their
  authors.}
  \label{tab:parity}
  \scriptsize
  \setlength{\tabcolsep}{6pt}
  \begin{tabular}{l cc}
    \toprule
    Configuration & $k{=}0$ & $k{=}10$ \\
    \midrule
    HALO (native Hz, per-sensor text)          & \PH & \PH \\
    HALO \emph{parity} (20\,Hz, neutral text)  & \PH & \PH \\
    \midrule
    Best baseline (data-matched group)         & \PH & \PH \\
    Best baseline (released-checkpoint group)  & \PH & \PH \\
    \bottomrule
  \end{tabular}
\end{table}

\section{Gravity as a Signed DC Feature}
\label{app:gravity}

Section~\ref{sec:sensor_grouping} keeps the quasi-static component as a signed per-axis feature
rather than normalising it away. Figure~\ref{fig:gravity} is the measurement behind that choice.

\begin{figure*}[t]
    \centering
    \setlength{\abovecaptionskip}{4pt}
  \resizebox{0.86\textwidth}{!}{% Retired on 2026-08-17. The previous posture similarities lacked a checked-in generating script,
% input-window manifest, and raw output. Reintroduce only after the signed-low-frequency ablation is
% reproduced and registered in EVIDENCE_LEDGER.md.
}
  \caption{\textbf{Why gravity is conditioned on rather than normalised away.} Three static
  postures, UCI-HAR waist phone, $200$ windows each through the real tokenizer path.
  \emph{Left:} the frequency path receives almost nothing (note the $\times10^{-3}$ axis) and
  separates no pair distinctly---all pairwise cosines fall within $0.84$--$0.88$. \emph{Right:}
  signed DC places lying nearly orthogonal to standing ($-0.04$), because gravity has
  redistributed across axes. The limit is shown too: standing and sitting stay at $0.99$ in DC,
  the same pair Section~\ref{sec:geometry} finds inseparable in text.}
  \label{fig:gravity}
\end{figure*}

\section{Tokenizer Guardrails}
\label{app:tokenizer}

\begin{figure}[h]
    \centering
    \setlength{\abovecaptionskip}{4pt}
  % =====================================================================
% fig_tokenizer.tex  --  the physical-Hz constant-Q filterbank
%
% DATA PROVENANCE: band centers are the exact values the implementation
% builds, c_k = f_min (f_max/f_min)^{k/(K-1)}, K=32, f_min=0.3, f_max=15 Hz
% (model/tokenizer/filterbank.py); sigma_k = c_k/(2Q) with Q=4.  The Nyquist
% guard lines are FB_NYQUIST_MARGIN * r/2 with margin 0.9.  Observable-band
% counts (26/32 at 20 Hz, 27/32 at 25 Hz, 32/32 at 50 and 100 Hz) were
% enumerated from the same constants on 2026-07-26.
% =====================================================================
\resizebox{\columnwidth}{!}{%
\begin{tikzpicture}[font=\scriptsize]

% ---------------- filterbank response ----------------
\begin{axis}[
  width=\columnwidth, height=3.5cm,
  xmode=log, log basis x=10,
  xmin=0.28, xmax=17, ymin=0, ymax=1.46,
  xtick={0.3,0.5,1,2,5,10,15},
  xticklabels={0.3,0.5,1,2,5,10,15},
  xticklabel style={font=\scriptsize},
  ytick=\empty,
  xlabel={physical frequency (Hz) --- the \emph{same} axis at every sampling rate},
  xlabel style={font=\scriptsize,yshift=2pt},
  axis y line=none, axis x line*=bottom,
  axis line style={draw=halogray!70},
  scale only axis, clip=false,
]
% --- a representative subset of the K=32 log-spaced constant-Q bands (Q=4) ---
\foreach \c in {0.300,0.497,0.823,1.364,2.259,3.743,6.201,10.272,15.000} {
  % the max(-25,.) clamp keeps pgfmath's exp() from underflowing to a
  % "Dimension too large" error in the far tails of the narrow low-Hz bands
  \edef\tmp{\noexpand\addplot[domain=0.28:17,samples=260,smooth,draw=haloblue!75,
      line width=0.5pt]{exp(max(-25,-((x-\c)^2)/(2*(\c/8)^2)))};}
  \tmp
}
% --- Nyquist guard lines ---
\addplot[dashed,line width=0.7pt,draw=haloorange] coordinates {(9.0,0) (9.0,1.18)};
\addplot[dashed,line width=0.7pt,draw=haloorange!60] coordinates {(11.25,0) (11.25,1.18)};
\node[font=\scriptsize,color=haloorange,anchor=south east,inner sep=1pt]
  at (axis cs:8.9,1.18) {20\,Hz guard};
\node[font=\scriptsize,color=haloorange!70!black,anchor=south west,inner sep=1pt]
  at (axis cs:11.4,1.18) {25\,Hz};
% --- masked-band shading ---
\addplot[draw=none,fill=haloorange!10,forget plot]
  coordinates {(9.0,0) (17,0) (17,1.15) (9.0,1.15)} \closedcycle;
\node[font=\scriptsize,color=haloorange!85!black,anchor=north,align=center]
  at (axis cs:12.4,0.80) {\textbf{masked}\\ at 20\,Hz};
\node[font=\scriptsize,color=haloblue!80!black,anchor=south west]
  at (axis cs:0.30,1.30) {$K{=}32$ log-spaced constant-$Q$ bands, $Q{=}4$};
\end{axis}

% ---------------- what the token actually carries ----------------
% Real mask vectors from PhysicalFilterbankTokenizer._observability_masks
% at the three (rate, patch-duration) combinations named on the left.
\begin{scope}[yshift=-1.55cm, xshift=1.55cm]
\node[font=\scriptsize\bfseries,anchor=west] at (-1.55,0.62)
  {each token carries the band energies \emph{and} what the configuration could measure:};
\node[fbblur] at (0.000,-0.000) {};
\node[fbblur] at (0.148,-0.000) {};
\node[fbblur] at (0.296,-0.000) {};
\node[fbblur] at (0.444,-0.000) {};
\node[fbblur] at (0.592,-0.000) {};
\node[fbblur] at (0.740,-0.000) {};
\node[fbblur] at (0.888,-0.000) {};
\node[fbok] at (1.036,-0.000) {};
\node[fbok] at (1.184,-0.000) {};
\node[fbok] at (1.332,-0.000) {};
\node[fbok] at (1.480,-0.000) {};
\node[fbok] at (1.628,-0.000) {};
\node[fbok] at (1.776,-0.000) {};
\node[fbok] at (1.924,-0.000) {};
\node[fbok] at (2.072,-0.000) {};
\node[fbok] at (2.220,-0.000) {};
\node[fbok] at (2.368,-0.000) {};
\node[fbok] at (2.516,-0.000) {};
\node[fbok] at (2.664,-0.000) {};
\node[fbok] at (2.812,-0.000) {};
\node[fbok] at (2.960,-0.000) {};
\node[fbok] at (3.108,-0.000) {};
\node[fbok] at (3.256,-0.000) {};
\node[fbok] at (3.404,-0.000) {};
\node[fbok] at (3.552,-0.000) {};
\node[fbok] at (3.700,-0.000) {};
\node[fbok] at (3.848,-0.000) {};
\node[fbok] at (3.996,-0.000) {};
\node[fbok] at (4.144,-0.000) {};
\node[fbok] at (4.292,-0.000) {};
\node[fbok] at (4.440,-0.000) {};
\node[fbok] at (4.588,-0.000) {};
\node[font=\scriptsize,anchor=east] at (-0.16,-0.000) {100\,Hz, 1.5\,s patch};
\node[font=\scriptsize,anchor=west,color=halogray] at (4.748,-0.000) {32/32 obs., 25/32 res.};
\node[fbblur] at (0.000,-0.550) {};
\node[fbblur] at (0.148,-0.550) {};
\node[fbblur] at (0.296,-0.550) {};
\node[fbblur] at (0.444,-0.550) {};
\node[fbblur] at (0.592,-0.550) {};
\node[fbblur] at (0.740,-0.550) {};
\node[fbblur] at (0.888,-0.550) {};
\node[fbok] at (1.036,-0.550) {};
\node[fbok] at (1.184,-0.550) {};
\node[fbok] at (1.332,-0.550) {};
\node[fbok] at (1.480,-0.550) {};
\node[fbok] at (1.628,-0.550) {};
\node[fbok] at (1.776,-0.550) {};
\node[fbok] at (1.924,-0.550) {};
\node[fbok] at (2.072,-0.550) {};
\node[fbok] at (2.220,-0.550) {};
\node[fbok] at (2.368,-0.550) {};
\node[fbok] at (2.516,-0.550) {};
\node[fbok] at (2.664,-0.550) {};
\node[fbok] at (2.812,-0.550) {};
\node[fbok] at (2.960,-0.550) {};
\node[fbok] at (3.108,-0.550) {};
\node[fbok] at (3.256,-0.550) {};
\node[fbok] at (3.404,-0.550) {};
\node[fbok] at (3.552,-0.550) {};
\node[fbok] at (3.700,-0.550) {};
\node[fbno] at (3.848,-0.550) {};
\node[fbno] at (3.996,-0.550) {};
\node[fbno] at (4.144,-0.550) {};
\node[fbno] at (4.292,-0.550) {};
\node[fbno] at (4.440,-0.550) {};
\node[fbno] at (4.588,-0.550) {};
\node[font=\scriptsize,anchor=east] at (-0.16,-0.550) {20\,Hz, 1.5\,s patch};
\node[font=\scriptsize,anchor=west,color=halogray] at (4.748,-0.550) {26/32 obs., 25/32 res.};
\node[fbblur] at (0.000,-1.100) {};
\node[fbblur] at (0.148,-1.100) {};
\node[fbblur] at (0.296,-1.100) {};
\node[fbblur] at (0.444,-1.100) {};
\node[fbblur] at (0.592,-1.100) {};
\node[fbblur] at (0.740,-1.100) {};
\node[fbblur] at (0.888,-1.100) {};
\node[fbblur] at (1.036,-1.100) {};
\node[fbblur] at (1.184,-1.100) {};
\node[fbblur] at (1.332,-1.100) {};
\node[fbblur] at (1.480,-1.100) {};
\node[fbblur] at (1.628,-1.100) {};
\node[fbblur] at (1.776,-1.100) {};
\node[fbblur] at (1.924,-1.100) {};
\node[fbblur] at (2.072,-1.100) {};
\node[fbblur] at (2.220,-1.100) {};
\node[fbok] at (2.368,-1.100) {};
\node[fbok] at (2.516,-1.100) {};
\node[fbok] at (2.664,-1.100) {};
\node[fbok] at (2.812,-1.100) {};
\node[fbok] at (2.960,-1.100) {};
\node[fbok] at (3.108,-1.100) {};
\node[fbok] at (3.256,-1.100) {};
\node[fbok] at (3.404,-1.100) {};
\node[fbok] at (3.552,-1.100) {};
\node[fbok] at (3.700,-1.100) {};
\node[fbok] at (3.848,-1.100) {};
\node[fbok] at (3.996,-1.100) {};
\node[fbok] at (4.144,-1.100) {};
\node[fbok] at (4.292,-1.100) {};
\node[fbok] at (4.440,-1.100) {};
\node[fbok] at (4.588,-1.100) {};
\node[font=\scriptsize,anchor=east] at (-0.16,-1.100) {100\,Hz, 0.5\,s patch};
\node[font=\scriptsize,anchor=west,color=halogray] at (4.748,-1.100) {32/32 obs., 16/32 res.};
\node[font=\scriptsize,anchor=west,align=left] at (-1.55,-1.62)
  {\tikz\node[fbok]{};\ measured \quad
   \tikz\node[fbblur]{};\ present but unresolved (flagged) \quad
   \tikz\node[fbno]{};\ beyond Nyquist (masked)};
\end{scope}
\end{tikzpicture}}


  \caption{\textbf{The filterbank, and what a token carries.} \emph{Top:} $K{=}32$ Gaussian bands
  log-spaced over $0.3$--$15$\,Hz; because a native-rate zero-padded rDFT maps bin $m$ to $mr/S$,
  one bank reads every rate without resampling. Dashed lines mark the Nyquist guard for the two
  lowest corpus rates. \emph{Bottom:} the real observability and resolution vectors emitted
  alongside the energies---part of the token, not diagnostics.}
  \label{fig:tokenizer}
\end{figure}

Zero-padding beyond the true patch length is frequency-domain interpolation only and does not
alter band energies; we verified that $S{=}256$ and $S{=}512$ give band-energy cosine $1.000000$
and identical masks at $20$, $25$, $50$, and $100$\,Hz. Patch lengths exceeding $S$ raise an error
rather than silently truncating.

\section{Phase-A Specification}
\label{app:phasea}

The EMA teacher is updated only after an optimizer step (decay $0.996$) and is checkpointed with
the model. Cross-placement positives are drawn under a quota so that the six-placement recordings
cannot dominate; positive identity comes from explicit physical-event identifiers written during
conversion and carried into the pretraining index, never inferred from equal array lengths or row
alignment, and never from an activity label. Per-objective gradient norms, VICReg component
values, minimum feature standard deviation, effective rank, and positive similarity by pair type
are logged every step; non-finite values or representation collapse fail the run rather than
being smoothed over.

\section{Phase-A Objective Ablation}
\label{app:objectives}

\begin{table}[h]
  \centering
  \caption{Phase-A objective ablation, macro-F1 over the held-out datasets. The last two rows are
  controls, not variants: they test whether the non-contrastive and label-free choices were worth
  making.}
  \label{tab:objective_ablation}
  \scriptsize
  \setlength{\tabcolsep}{3pt}
  \begin{tabular}{l cc}
    \toprule
    Phase-A configuration & $k{=}0$ & $k{=}10$ \\
    \midrule
    Full (A1 fixed$+$EMA, VICReg, TF, placement, rail) & \PH & \PH \\
    \quad$-$ EMA-latent target (fixed A1 only)      & \PH & \PH \\
    \quad$-$ augmentation VICReg                    & \PH & \PH \\
    \quad$-$ time--frequency consistency            & \PH & \PH \\
    \quad$-$ cross-placement positives              & \PH & \PH \\
    \quad$-$ physical grounding rail                & \PH & \PH \\
    \midrule
    VICReg $\rightarrow$ NT-Xent (contrastive control) & \PH & \PH \\
    Label-supervised recipe (SupCon; \emph{uses labels}) & \PH & \PH \\
    \bottomrule
  \end{tabular}
\end{table}

\section{Phase-B Episodic Training}
\label{app:phaseb}

The decoder is trained on episodes that vary, independently: the candidate budget
($10$--$100\%$ of the vocabulary); the number of true-label memory exemplars
($0$, $1$, $2$, $4$, $8$, or all); the support available for competing labels; whether memory is
same-configuration, cross-configuration only, or missing the query's configuration entirely; the
kind of distractors present (random, language-near, motion-family-near, physically confusable);
and whether the truth is present at all. Complete label \emph{families} are held out of decoder
training and reserved for model selection, so candidate-text reasoning is always tested on
concepts the decoder never fit. Head collapse in the retrieval projections is resisted by output
decorrelation, projection diversity, and contribution-usage regularisation.
